% Guia do professor — Capítulo 12: Massas de Ar e Frentes
\section{Capítulo 12 --- Massas de Ar e Frentes}

\subsection*{Visão geral e ritmo}

A ponte entre a circulação geral (Cap.~11) e os ciclones (Cap.~13): onde
o ar ganha identidade (massas) e onde as identidades colidem (frentes).
Capítulo muito visual — cortes verticais e meteogramas rendem mais que
equações, com duas exceções nobres: Margules e a frontogênese
exponencial. \textbf{Ritmo: 2 aulas de 2\,h.}

\begin{itemize}
  \item \textbf{Aula 1} — anticiclones como fábricas, classificação de
        Bergeron com $(\theta_w, r)$, gênese fria (EDO no quadro),
        modificação em trânsito. Caso~1 (10\,min).
  \item \textbf{Aula 2} — anatomia comparada fria × quente, Margules,
        frontogênese por deformação, oclusões e linha seca. Casos~2--4
        (20\,min; o meteograma do Caso~4 fecha bem a aula).
\end{itemize}

\subsection*{Objetivos de aprendizagem}

\begin{enumerate}
  \item classificar massas de ar por origem e pelo par conservado
        $(\theta_w, r)$;
  \item estimar o tempo de gênese fria pela EDO radiativa linearizada;
  \item calcular a inclinação frontal por Margules e ligá-la à
        nebulosidade de cada frente;
  \item explicar a frontogênese por deformação e sua taxa exponencial;
  \item reconhecer passagens frontais em meteogramas e cartas
        (isóbaras em ``V'', giro de vento no HS).
\end{enumerate}

\subsection*{Pré-requisitos a reativar}

$\theta_w$ e $r$ como etiquetas conservadas (Cap.~4); fluxos bulk
(Cap.~3); geostrofia e hidrostática (Caps.~1, 10); Stefan--Boltzmann
(Cap.~2); leitura de cartas (Cap.~9).

\subsection*{Armadilhas conceituais frequentes}

\begin{itemize}
  \item \textbf{``Frente é uma linha''}: é uma \emph{zona} 3-D com
        $\sim$100\,km de largura e uma cunha quase deitada (1/100!) —
        o desenho em escala real é chocante e vale fazer.
  \item \textbf{Giro de vento no HS}: N$\to$SW na fria (não S$\to$NW,
        que é o espelho do HN dos livros importados). Insistir com a
        rosa dos ventos no quadro.
  \item \textbf{``A frente fria é fria''}: o que define é o
        \emph{avanço} do ar frio; num pós-frontal de céu limpo a tarde
        pode ser mais quente que a manhã pré-frontal.
  \item \textbf{Margules sem cisalhamento}: sem $\Delta U$ ao longo da
        frente não há interface inclinada estacionária — frentes são
        estruturas dinâmicas, não ``paredes de densidade''.
  \item \textbf{Oclusão comparando $T$}: compara-se $\theta_w$ das
        massas frias (T5); com $T$ crua o critério erra sob inversões.
\end{itemize}

\subsection*{Demonstração de baixo custo}

Tanque de água com divisória: água fria tingida de um lado, quente do
outro; ao remover a divisória a cunha fria avança e a inclinação
aparece (sem rotação ela desaba — deixe explícito que é Margules SEM o
$f$, por isso não estaciona). Meteograma real: baixar a série da
estação do IAG/aeroporto na última friagem e sobrepor ao Caso~4.

\subsection*{Problemas sugeridos do livro (exercícios do Cap.~12)}

Aquecimento: classificação de massas para cidades dadas; tempo de
gênese com números diferentes. Aprofundamento: inclinação frontal com
dados de sondagem dupla; frontogênese com campo de deformação dado.
Síntese: ``por que o HS tem frentes mais zonais e menos violentas que o
HN?''. \emph{Confira a numeração na sua cópia (v1.02b).}

\subsection*{Uso do notebook \texttt{cap12\_frentes.ipynb}}

\begin{center}
\begin{tabular}{lll}
\hline
Caso & Conteúdo & Momento sugerido \\
\hline
1 & gênese cP: EDO radiativa ($T^4$) vs.\ linearizada & Aula 1 \\
2 & inclinação de Margules: varredura $\Delta T$, $\Delta U$ & Aula 2 \\
3 & frontogênese por deformação (mini-NWP, exata) & Aula 2 \\
4 & meteograma sintético da frente fria & fecho da Aula 2 \\
\hline
\end{tabular}
\end{center}

O Caso~3 é deliberadamente a ``previsão numérica em miniatura'' da vez:
advecção com solução exata — compare com a versão em grade (upwind) do
Cap.~3 e anuncie o Cap.~20.

\subsection*{Exercícios complementares e soluções}

Nas notas: T1--T5 e N1--N4. Soluções em \texttt{cap12\_solucoes.ipynb}
(\emph{não distribuir}): T1 e T2 verificados com \texttt{sympy}; N3
responde uma pergunta genuína (rotação não frontogeniza). Pares:
T2$\leftrightarrow$N1, T1$\leftrightarrow$N2, T3$\leftrightarrow$N3.

\subsection*{Conexões com o resto do curso}

Anticiclones e subsidência (Caps.~10--11); $(\theta_w, r)$ (Cap.~4);
nuvens frontais (Cap.~6); meteograma e cartas (Cap.~9); a frente que
ondula e fecha em ciclone é o Cap.~13 inteiro; linha seca e severo
(Cap.~14); ZCAS e chuvas persistentes (Cap.~21, variabilidade).
